\documentclass[11pt,a4paper]{article}

%,fonts ------------------------------------------------------------
\usepackage[T1]{fontenc}
\usepackage[utf8]{inputenc}
\usepackage{mathpazo}                           % Palatino body + math
\usepackage[scaled=0.92]{helvet}            % Helvetica for sans
\linespread{1.07}

%,layout -----------------------------------------------------------
\usepackage[a4paper,
            top=2.4cm, bottom=2.4cm,
            left=2.6cm, right=2.6cm,
            headheight=14pt]{geometry}
\usepackage[protrusion=true,expansion=false]{microtype}
\usepackage{parskip}
\usepackage{xcolor}
\usepackage{amsmath}
\usepackage{enumitem}
\setlist{leftmargin=*,labelsep=0.6em,itemsep=2pt,topsep=4pt}

%,accent colour: muted navy / slate -------------------------------
\definecolor{ink}{HTML}{1F2A37}
\definecolor{slate}{HTML}{2D4356}
\definecolor{rule}{HTML}{C5CBD3}
\definecolor{codebg}{HTML}{F6F6F4}

%,headings: sans, with a thin rule below the section --------------
\usepackage{titlesec}
\titleformat{\section}
  {\normalfont\sffamily\large\bfseries\color{slate}}
  {\textcolor{slate}{\thesection \enskip\textbar\enskip}}{0pt}{}
  [\vspace{-2pt}{\color{rule}\hrule height 0.4pt}\vspace{2pt}]
\titlespacing*{\section}{0pt}{18pt}{6pt}

\titleformat{\subsection}
  {\normalfont\sffamily\normalsize\bfseries\color{slate}}{\thesubsection}{0.6em}{}
\titlespacing*{\subsection}{0pt}{10pt}{2pt}

\titleformat{\paragraph}[runin]
  {\normalfont\sffamily\small\bfseries\color{ink}}{}{0pt}{}[\enspace]
\titlespacing*{\paragraph}{0pt}{6pt}{0.4em}

%,header / footer -------------------------------------------------
\usepackage{fancyhdr}
\pagestyle{fancy}
\fancyhf{}
\renewcommand{\headrulewidth}{0pt}
\renewcommand{\footrulewidth}{0pt}
\fancyhead[L]{\footnotesize\sffamily\color{slate}IoT Course — Exercise 2}
\fancyhead[R]{\footnotesize\sffamily\color{slate}BLE denture localization}
\fancyfoot[C]{\footnotesize\sffamily\color{slate}\thepage}

%,tables: booktabs only, no colour blocks -------------------------
\usepackage{booktabs}
\usepackage{tabularx}
\usepackage{array}
\newcolumntype{Y}{>{\raggedright\arraybackslash}X}
\renewcommand{\arraystretch}{1.18}

%,code listings: minimal, thin left bar ---------------------------
\usepackage{listings}
\lstdefinestyle{pseudo}{
    backgroundcolor=\color{codebg},
    basicstyle=\ttfamily\footnotesize\color{ink},
    breaklines=true,
    frame=leftline,
    framerule=1.4pt,
    rulecolor=\color{slate},
    framesep=8pt,
    xleftmargin=14pt,
    xrightmargin=2pt,
    aboveskip=8pt,
    belowskip=6pt,
    keywordstyle=\bfseries\color{slate},
    commentstyle=\itshape\color{slate!60!black},
    morekeywords={init,loop,forever,if,elif,else,while,for,in,function,return,
                  and,or,not,wait_for,sleep,send_BLE_beacon,configure,
                  enable_interrupt,read,now}
}

\usepackage[hidelinks]{hyperref}

%,title block -----------------------------------------------------
\title{IoT - Challenge 4 Homework - Part 2}
\author{Raul Agolli - Team Leader \\\\ Gabriele Lorenzo Singe}
\date{May 2026}

\begin{document}
\maketitle
\thispagestyle{fancy}

%,context ---------------------------------------------------------
\noindent\textsc{Problem.}\enskip
The original firmware advertises every $10\,$s regardless of context.
For most of the day the denture is either \emph{in the patient's mouth}
(localization is not needed) or \emph{sitting still in the same place}
(repeated beacons add no information). The course's hardware lecture is
explicit on the point: radio energy dominates the budget, idle TX should
be avoided, and the radio must be switched off as soon as possible. The
fix is to replace periodic beaconing with an \emph{event-driven} scheme
driven by an accelerometer interrupt, with a slow heartbeat as a
keep-alive.

% =====================================================================
\section{Beaconing logic}
% =====================================================================

A three-state machine drives the radio.

\begin{tabularx}{\linewidth}{@{} l Y l @{}}
\toprule
\textsc{state} & \textsc{when}\ /\ \textsc{what} & \textsc{rate} \\
\midrule
\textsf{DEEP-SLEEP}
  & No motion for $>30\,$s, or denture worn. Radio OFF; MCU in deep
    sleep; accelerometer in low-power motion-detect mode. Slow
    heartbeat for keep-alive.
  & 1 / 5\,min \\
\addlinespace[2pt]
\textsf{ACTIVE}
  & Accelerometer wake-on-motion interrupt fired. Fast advertising for
    dense RSSI sampling by the gateways.
  & 1 / 1--2\,s \\
\addlinespace[2pt]
\textsf{SETTLING}
  & Motion has stopped for $>5\,$s but $<30\,$s. Short burst of final
    beacons so the resting position is well sampled, then back to
    \textsf{DEEP-SLEEP}.
  & 3 over 5\,s \\
\bottomrule
\end{tabularx}

\paragraph{Key insight.}
Localization is only useful while the denture has \emph{just moved} —
the caregiver needs to know where the patient set it down. When the
denture is static the gateways already have a valid fix.

% =====================================================================
\section{Additional sensors}
% =====================================================================

\begin{description}[font=\sffamily\bfseries\color{slate}]
\item[Accelerometer (mandatory).]
A 3-axis MEMS accelerometer with hardware wake-on-motion interrupt.
The sensor draws a few $\mu$A in motion-detect mode and asserts an
interrupt line straight to the BLE SoC — the exact wake-up pattern
described in the hardware-abstraction lecture (event-driven instead of
polling).

\item[Body-contact thermistor (optional, recommended).]
When the temperature is close to $36\,^\circ$C the denture is in the
mouth and localization can be \emph{suppressed entirely} (the patient
has it). This eliminates beaconing during the many hours per day the
denture is worn.
\end{description}

% =====================================================================
\section{Improved firmware}
% =====================================================================

\begin{lstlisting}[style=pseudo]
init:
    accel.configure(threshold = 0.1 g, ODR = 25 Hz, mode = MOTION_DETECT)
    accel.enable_interrupt(pin = INT1)        # wakes the MCU
    temp.configure(low_power_mode)
    send_BLE_beacon()                         # first position fix
    last_motion = now()
    state = DEEP_SLEEP

loop forever:

    if state == DEEP_SLEEP:
        # MCU + radio OFF; only accelerometer and RTC still running
        wait_for(motion_interrupt OR rtc_timeout(5 min))
        if motion_interrupt and not in_mouth():
            state = ACTIVE
            last_motion = now()
        else:                                 # 5-min heartbeat
            if not in_mouth():
                send_BLE_beacon()

    elif state == ACTIVE:
        send_BLE_beacon()                     # fast advertising
        sleep(1 s)
        if motion_interrupt_fired_recently():
            last_motion = now()
        if (now() - last_motion) > 5 s:
            state = SETTLING

    elif state == SETTLING:
        # small burst so the resting position is well sampled
        for i in 1..3:
            send_BLE_beacon()
            sleep(2 s)
        state = DEEP_SLEEP

function in_mouth():
    T = temp.read()
    return T > 30 C        # oral T ~36 C; off-mouth drops below 30 C
\end{lstlisting}

% =====================================================================
\newpage
\section{Commercial components}
% =====================================================================

\begin{tabularx}{\linewidth}{@{} l l Y @{}}
\toprule
\textsc{function} & \textsc{part} & \textsc{key spec} \\
\midrule
BLE 5 SoC                & Nordic \textsf{nRF52810} / \textsf{nRF52832}
& ARM Cortex-M4 @ $64\,$MHz; $1.5\,\mu$A System-ON sleep with full
  RAM retention; $\approx 4.6\,$mA TX at $0\,$dBm
  \emph{(Nordic product brief)}. \\
\addlinespace[2pt]
Wake-on-motion accel.\   & ST \textsf{LIS2DH12} \emph{or} Bosch \textsf{BMA400}
& 3-axis MEMS; LIS2DH12 $\approx 2\,\mu$A at $10\,$Hz ODR in
  low-power mode; BMA400 $<1\,\mu$A in self-wake-up; dedicated INT
  pin. \\
\addlinespace[2pt]
Temperature sensor       & TI \textsf{TMP117} \emph{or} ST \textsf{STTS22H}
& I$^2$C; TMP117 is designed for clinical thermometers (ASTM E1112,
  $\pm 0.1\,^\circ$C); STTS22H consumes $1.75\,\mu$A in one-shot mode. \\
\addlinespace[2pt]
Battery                  & \textsf{CR2032} coin cell
& $\approx 220\,$mAh at $3\,$V (mid of the manufacturer
  $210$--$240\,$mAh range). \\
\bottomrule
\end{tabularx}

% =====================================================================
\section{Energy-saving estimate}
% =====================================================================

\subsection*{Working numbers}

\begin{description}[font=\sffamily\small\bfseries\color{slate}]
\item[Charge per advertising event.]
A full event on the three advertising channels lasts $\approx 5\,$ms
with an average current $\approx 2.85\,$mA (peak $\approx 5\,$mA on
the actual TX bursts). That gives $\approx 14\,\mu$C $\approx
\mathbf{4\,\text{nAh}}$ per beacon, the figure produced by Nordic's
\emph{nRF52 online power profiler}, which we use as representative.
\item[Idle current.]
$1.5\,\mu$A (nRF52810 System-ON, RAM ret.) $+\ 1.5\,\mu$A
(LIS2DH12 low-power, $10\,$Hz) $\approx 3\,\mu$A
$\Rightarrow 72\,\mu$Ah/day. Identical in both firmwares, cancels in
the radio-saving ratio.
\item[Battery.] CR2032, $\approx 220\,$mAh at $3\,$V.
\item[Usage model.]
The denture is moved $\approx 15\,$min/day in total (taken out,
cleaned, repositioned); the other $\approx 23.75\,$h it sits still.
With the in-mouth sensor, the heartbeat is further suppressed during the
$\approx 20\,$h/day the denture is worn.
\end{description}

\subsection*{Beacons per day}

\begin{tabularx}{\linewidth}{@{} Y r r @{}}
\toprule
                                        & \textsc{original} & \textsc{improved} \\
\midrule
Periodic beacons (1 every $10\,$s, $24\,$h) & $24\cdot 3600/10 = 8\,640$ &,\\
Motion-driven ($15\,$min, $1\,$beacon/s)    &,     & $900$            \\
Heartbeat ($23.75\,$h, 1 per $5\,$min)      &,     & $\approx 285$    \\
Settling burst ($\sim 10$ episodes, 3 each) &,     & $\approx 30$     \\
\midrule
\textbf{Total}                              & \textbf{8\,640} & $\mathbf{\approx 1\,215}$ \\
\bottomrule
\end{tabularx}

\smallskip
\noindent With the in-mouth temperature sensor active, heartbeats are
emitted only during the $\approx 4\,$h/day the denture is out of
mouth, so heartbeats drop to $\approx 48$ and the total becomes
$\approx \mathbf{978}$ beacons/day.

\subsection*{Radio-energy comparison}

\begin{tabularx}{\linewidth}{@{} Y r r @{}}
\toprule
\textsc{scenario}                       & \textsc{beacons/day} & \textsc{TX charge/day} \\
\midrule
Original (periodic $10\,$s)             & $8\,640$ & $\approx 34.6\,\mu$Ah \\
Improved (no temperature sensor)        & $1\,215$ & $\approx 4.9\,\mu$Ah  \\
Improved $+$ in-mouth detection         & $978$    & $\approx 3.9\,\mu$Ah  \\
\bottomrule
\end{tabularx}

\smallskip
\noindent
Saving on the radio component:\par
\hspace*{1em}--- improved, no temperature sensor:\quad
$1 - 4.9/34.6 \approx \mathbf{86\,\%}$\par
\hspace*{1em}--- improved + in-mouth detection:\quad
$1 - 3.9/34.6 \approx \mathbf{89\,\%}$

\subsection*{Battery lifetime}

Adding the identical idle current of $72\,\mu$Ah/day:

\begin{tabularx}{\linewidth}{@{} Y r r @{}}
\toprule
\textsc{scenario}                       & \textsc{total charge/day} & \textsc{CR2032 lifetime}$^*$ \\
\midrule
Original                                & $72 + 34.6 \approx 107\,\mu$Ah  & $\approx 2060\,$d \,($\approx 5.6\,$y) \\
Improved (no temp)                      & $72 + 4.9 \approx 77\,\mu$Ah    & $\approx 2860\,$d \,($\approx 7.8\,$y) \\
Improved $+$ in-mouth                   & $72 + 3.9 \approx 76\,\mu$Ah    & $\approx 2900\,$d \,($\approx 7.9\,$y) \\
\bottomrule
\end{tabularx}

\smallskip

\paragraph{Result}
Radio activity is cut by $86$--$89\,\%$. Because the idle current then
dominates the residual budget, the overall daily energy drops by
$\approx 28\,\%$ and battery lifetime is extended by roughly
$\mathbf{1.4\times}$. Critically, the localization \emph{quality} also
improves: while the denture is actually moving, gateways now receive one
beacon every second instead of one every ten seconds, so the RSSI
samples used for the fix are denser exactly when the answer matters.

\end{document}